% ===================================================================== Chapter 16
\chapter{Uncertainty and Limitations}\label{ch:unc}
Nine sources of uncertainty were identified and are kept separate (Table~\ref{tab:unc}). They are of different kinds -- discretisation
errors, model-form and idealisation errors, data scatter and assumptions -- several have a known sign and some are not quantified. A combined
uncertainty percentage would therefore have no mathematical basis and is not formed. The evidence supports the ranking primary (end restraint
and stability modelling) $\gg$ secondary (thermal field and material data) $\gg$ smaller (mapping, structural mesh, Poisson's ratio).

%% generated by gen_tables.py - RE-ANALYSIS 2026
\begin{table}[htbp]
\centering
\footnotesize
\caption{Uncertainty register: nine separate sources (no combined percentage is formed)}
\label{tab:unc}
\begin{tabular}{l>{\raggedright\arraybackslash}p{0.20\linewidth}>{\raggedright\arraybackslash}p{0.50\linewidth}>{\raggedright\arraybackslash}p{0.11\linewidth}}
\toprule
\# & Source & Size / evidence & Class \\
\midrule
7 & End-restraint definition & $\lambda_1$ = 1.108 / 2.232 / 4.300 (S1 / S3 / S2) at identical static stress: factor 3.9; mechanism switches & primary \\
8 & Geometric imperfection & not modelled in the Mechanical study (numerical amplitudes only in the additional elastic LS-DYNA analysis, Section~\ref{sec:lsdyna}); reduces real capacity below $\lambda_1$ by an unknown amount & primary \\
9 & Inelastic buckling & intermediate columns ($KL/r$ 49.7--58.3 $<$ $C_c \approx$ 60); Johnson estimate 1.058 for S1 & primary \\
2 & CFD model form & not quantified; 19.1 K CFD--analytical gap on $T_{max}$; $\approx$30 K on the mean rise; $\lambda_1 \approx$ 0.98 on the Section 2 temperature basis & secondary \\
1 & CFD mesh & $T_{max}$ $-$5.1 / +1.6 K; LC2 stress $-$10.1 / +2.8 MPa; $\lambda_1$ +1.8 / $-$0.5\,\% & secondary \\
5 & Material data & about $\pm$2\,\% on $E$ and $\alpha$; $S_y$ typical values & secondary \\
3 & Temperature mapping & $\le$ 0.081 K interior; LC1 inlet peak bound $\pm$9.8 MPa & smaller \\
4 & Structural mesh & LC2 peak $\le 4.9\times10^{-4}$, $\lambda_1 \le 5.4\times10^{-5}$; LC1 bore $-$2.1\,\% & smaller \\
6 & Poisson's ratio & about 1.4\,\% on LC1; none on LC2 & smaller \\
\bottomrule
\end{tabular}
\par\smallskip\parbox{\linewidth}{\footnotesize\raggedright Source: \texttt{11\_Final\_Audit/MASTER\_UNCERTAINTIES.md}. Sources are of different kind and sign (discretisation, model form, data scatter, assumption); a root-sum-square or linear sum would have no mathematical basis.}
\end{table}


\section{Data-Availability Limitation}
The original internship files were not retained, so no result can be compared with the analysis made during the internship, and nothing
in this report is a recovered value. The analysed problem is representative of the internship problem, not a reproduction of it.

\section{CFD Model Form}
The turbulence and property models were not varied. Indicators of the model-form uncertainty are the 9\,\% difference between the CFD Nusselt
number and the property-corrected correlation, the 10.4\,\% difference in friction factor (mostly a gas-heating effect, about five points
unexplained), the 19.1~K difference in maximum solid temperature from the analytical model and the difference of about 30~K in mean
temperature rise. The last one matters structurally: on the Section 2 temperature basis the S1 buckling factor would be about 0.98 instead of
1.108. Uniform heating, neglected radiation and the assumed inlet turbulence intensity add further, unquantified, model-form uncertainty.

\section{CFD Mesh}
The medium mesh is approximately convergent. Its band against the extrapolated solution is $-$5.1 / +1.6~K on the maximum solid temperature, which
propagates to $-$10.1 / +2.8~MPa on the LC2 stress, $-$13.7 / +4.3~MPa on the local peak, +1.8 / $-$0.5\,\% on $\lambda_1$ and $-$32 / +9~\textmu m on
the free growth. The medium field is slightly hot, so the structural results lean conservative.

\section{Temperature Mapping}
The transfer error is at most 0.081~K inside the source mesh. Within about 11~mm of the inlet the source node values carry the Fluent
reconstruction limitation of up to 2.6~K, which bounds the LC1 inlet peak at $\pm$9.8~MPa and the LC2 peak at 1.6\,\%. The re-numbering of source
nodes between otherwise identical cases changes the mapped field by at most 0.081~K and the LC2 peak by $9.5\times10^{-6}$.

\section{Structural Mesh}
The governing LC2 results and $\lambda_1$ change by at most $4.9\times10^{-4}$ and $5.4\times10^{-5}$ across seven meshes. The LC1 mid-span bore
stress is under-predicted by 2.1\,\% on mesh B (sign known). On one thickness check mesh the LC2 inlet-edge peak changed by $1.4\times10^{-4}$,
above the $10^{-4}$ class-A tolerance at a single node; this was reported, not relaxed.

\section{Material Data}
The properties are generic datasheet values, not lot-specific. An uncertainty of about $\pm$2\,\% on $E$ and $\alpha$ acts directly on the
restrained stress and end force ($\propto E\alpha$). The yield strengths are typical rather than minimum-guaranteed values, so the utilisation
is not a certified margin. No temperature-dependent stress--strain curve is available.

\section{Poisson's Ratio}
$\nu$ = 0.294 is assumed; the literature range is 0.284--0.294. It affects the LC1 radial and hoop stresses by about 1.4\,\% and the uniaxial LC2
axial stress not at all; $\lambda_1$ is affected only through the shear correction.

\section{End Restraint}
The real flange stiffness is undefined. The idealised supports change $\lambda_1$ from 1.108 (S1) to 2.232 (S3) and 4.300 (S2) at identical
static stress and switch the governing mechanism from bifurcation-first to first-yield-first. Rotational flexibility would lower $\lambda_1$ below
S1. This is the primary uncertainty of every stability statement in the report.

\section{Geometric Imperfections}
The tube is modelled perfectly straight, without out-of-straightness, eccentricity of the end load or residual stress. A real tube deflects
laterally below the bifurcation load and its capacity is lower; the size of the reduction is unknown without an imperfection-sensitive
nonlinear analysis. The additional LS-DYNA analysis (Section~\ref{sec:lsdyna}) quantifies the elastic response for three numerical imperfection
amplitudes; without measured imperfection data and an inelastic material model, the reduction for a real tube remains unknown.

\section{Inelastic Buckling}
All cases are intermediate columns ($KL/r$ = 49.7--58.3 under S1, below $C_c \approx$ 60) with an elastic buckling stress above the conventional
proportional limit. The linear elastic eigenvalue analysis cannot capture the resulting inelastic reduction; the conventional Johnson factors
(1.058 guided, 1.414 fixed--pinned, 1.581 fixed--fixed) indicate its direction but are not material-specific.

\section{Missing Experimental Validation}
No measured data exist. The CFD is verified against conservation laws, correlations and mesh studies, and the structural model against
closed-form solutions and equilibrium checks, but neither is validated against experiment. Agreement with correlations shows consistency,
not accuracy with respect to a physical duct.

\section{Interaction Effects}
The parametric study is one factor at a time over a narrow envelope ($\pm$10\,\% in velocity and heat flux, $\pm$20\,\% in thickness). Interactions
of about 8--11\,\% of the main effects were estimated for velocity and heat flux but not simulated. Transient start-up, thermal cycling and
fatigue are outside the scope.
